\newglossaryentry{event}{
    name={Уведомительное событие},
    description={зарегистрированное намерение выполнить доставку сообщения заданной аудитории с определённым шаблоном и полезной нагрузкой}
}

\newglossaryentry{dispatch}{
    name={Запуск доставки},
    description={экземпляр исполнения уведомительного события с фиксированной аудиторией и параметрами отправки}
}

\newglossaryentry{delivery}{
    name={Канальная доставка},
    description={обработка уведомления для конкретного получателя и конкретного канала связи в рамках одного запуска доставки}
}

\newglossaryentry{delivery_history}{
    name={История доставки},
    description={агрегированное представление статусов, позволяющее восстановить путь обработки от входной команды до итогового результата}
}

\newglossaryentry{template_term}{
    name={Шаблон},
    description={версионируемое описание содержимого уведомления для конкретного канала связи}
}

\newglossaryentry{provider}{
    name={Канальный шлюз},
    description={внешний или внутренний компонент, выполняющий фактическую передачу уведомления по выбранному каналу}
}

\newglossaryentry{outbox_pattern}{
    name={transactional outbox},
    description={паттерн надёжной публикации интеграционных сообщений, при котором бизнес-изменение и запись сообщения фиксируются в одной транзакции базы данных}
}

\newglossaryentry{idempotency}{
    name={идемпотентность},
    description={свойство операции, при котором повторный запуск с теми же входными параметрами не приводит к дополнительным побочным эффектам}
}

\newglossaryentry{deduplication}{
    name={дедупликация},
    description={обнаружение и подавление дублирующихся сообщений или операций по устойчивому ключу}
}

\newglossaryentry{redis_deduplication}{
    name={дедупликация через Redis},
    description={техническое подавление повторной обработки sender-команд через атомарную запись TTL-ключа в Redis}
}

\newglossaryentry{retry_term}{
    name={повторная попытка},
    description={контролируемое повторное выполнение операции после временной ошибки}
}

\newglossaryentry{fallback_routing}{
    name={fallback-маршрутизация},
    description={переназначение доставки в альтернативный канал связи после окончательной ошибки или исчерпания повторов в исходном канале}
}

\newglossaryentry{cancellation_mark}{
    name={отметка отмены},
    description={TTL-ключ в Redis, фиксирующий отмену dispatch или delivery и используемый sender-сервисами перед фактической отправкой}
}

\newglossaryentry{history_writer}{
    name={history-writer},
    description={сервис, который потребляет статусные события доставки и записывает append-only историю в ClickHouse}
}

\newglossaryentry{eventual_consistency}{
    name={eventual consistency},
    description={модель согласованности распределённой системы, в которой синхронизация состояний между сервисами происходит за конечное, но не мгновенное время}
}

\newglossaryentry{at_least_once}{
    name={at-least-once delivery},
    description={модель доставки сообщений, при которой каждое сообщение доставляется получателю не менее одного раза; возможны дубликаты, которые должны обрабатываться на уровне приложения}
}

\newglossaryentry{fault_tolerance}{
    name={fault tolerance},
    description={свойство системы сохранять работоспособность при частичных отказах компонентов}
}

\newglossaryentry{document_model}{
    name={документоориентированная модель данных},
    description={модель хранения, в которой связанные атрибуты предметной сущности размещаются в одном документе и обрабатываются как единое целое}
}
